%Version 3.1 December 2024
% See section 11 of the User Manual for version history
%
%%%%%%%%%%%%%%%%%%%%%%%%%%%%%%%%%%%%%%%%%%%%%%%%%%%%%%%%%%%%%%%%%%%%%%
%%                                                                 %%
%% Please do not use \input{...} to include other tex files.       %%
%% Submit your LaTeX manuscript as one .tex document.              %%
%%                                                                 %%
%% All additional figures and files should be attached             %%
%% separately and not embedded in the \TeX\ document itself.       %%
%%                                                                 %%
%%%%%%%%%%%%%%%%%%%%%%%%%%%%%%%%%%%%%%%%%%%%%%%%%%%%%%%%%%%%%%%%%%%%%

%%\documentclass[referee,sn-basic]{sn-jnl}% referee option is meant for double line spacing

%%=======================================================%%
%% to print line numbers in the margin use lineno option %%
%%=======================================================%%

%%\documentclass[lineno,pdflatex,sn-basic]{sn-jnl}% Basic Springer Nature Reference Style/Chemistry Reference Style

%%=========================================================================================%%
%% the documentclass is set to pdflatex as default. You can delete it if not appropriate.  %%
%%=========================================================================================%%

%%\documentclass[sn-basic]{sn-jnl}% Basic Springer Nature Reference Style/Chemistry Reference Style

%%Note: the following reference styles support Namedate and Numbered referencing. By default the style follows the most common style. To switch between the options you can add or remove “Numbered” in the optional parenthesis. 
%%The option is available for: sn-basic.bst, sn-chicago.bst%  
 
%%\documentclass[pdflatex,sn-nature]{sn-jnl}% Style for submissions to Nature Portfolio journals
%%\documentclass[pdflatex,sn-basic]{sn-jnl}% Basic Springer Nature Reference Style/Chemistry Reference Style
\documentclass[pdflatex,sn-mathphys-num]{sn-jnl}% Math and Physical Sciences Numbered Reference Style
%%\documentclass[pdflatex,sn-mathphys-ay]{sn-jnl}% Math and Physical Sciences Author Year Reference Style
%%\documentclass[pdflatex,sn-aps]{sn-jnl}% American Physical Society (APS) Reference Style
%%\documentclass[pdflatex,sn-vancouver-num]{sn-jnl}% Vancouver Numbered Reference Style
%%\documentclass[pdflatex,sn-vancouver-ay]{sn-jnl}% Vancouver Author Year Reference Style
%%\documentclass[pdflatex,sn-apa]{sn-jnl}% APA Reference Style
%%\documentclass[pdflatex,sn-chicago]{sn-jnl}% Chicago-based Humanities Reference Style

%%%% Standard Packages
%%<additional latex packages if required can be included here>

\usepackage{graphicx}%
\usepackage{multirow}%
\usepackage{amsmath,amssymb,amsfonts}%
\usepackage{amsthm}%
\usepackage{mathrsfs}%
\usepackage[title]{appendix}%
\usepackage{xcolor}%
\usepackage{textcomp}%
\usepackage{manyfoot}%
\usepackage{booktabs}%
\usepackage{algorithm}%
\usepackage{algorithmicx}%
\usepackage{algpseudocode}%
\usepackage{listings}%
\usepackage{placeins}

%%%%

%%%%%=============================================================================%%%%
%%%%  Remarks: This template is provided to aid authors with the preparation
%%%%  of original research articles intended for submission to journals published 
%%%%  by Springer Nature. The guidance has been prepared in partnership with 
%%%%  production teams to conform to Springer Nature technical requirements. 
%%%%  Editorial and presentation requirements differ among journal portfolios and 
%%%%  research disciplines. You may find sections in this template are irrelevant 
%%%%  to your work and are empowered to omit any such section if allowed by the 
%%%%  journal you intend to submit to. The submission guidelines and policies 
%%%%  of the journal take precedence. A detailed User Manual is available in the 
%%%%  template package for technical guidance.
%%%%%=============================================================================%%%%

%% as per the requirement new theorem styles can be included as shown below
\theoremstyle{thmstyleone}%
\newtheorem{theorem}{Theorem}%  meant for continuous numbers
%%\newtheorem{theorem}{Theorem}[section]% meant for sectionwise numbers
%% optional argument [theorem] produces theorem numbering sequence instead of independent numbers for Proposition
\newtheorem{proposition}[theorem]{Proposition}% 
%%\newtheorem{proposition}{Proposition}% to get separate numbers for theorem and proposition etc.

\theoremstyle{thmstyletwo}%
\newtheorem{example}{Example}%
\newtheorem{remark}{Remark}%

\theoremstyle{thmstylethree}%
\newtheorem{definition}{Definition}%

\raggedbottom
%%\unnumbered% uncomment this for unnumbered level heads

\begin{document}

\title[EEG Motor Action and Motor Imagery Classification]{EEG-Based Motor Action and Motor Imagery Classification Using Machine Learning}

\author[1]{\fnm{Harsada} \sur{K}}\email{cb.ai.u4aid24019@cb.students.amrita.edu}
\author[1]{\fnm{V} \sur{Tanay Darshan}}\email{cb.ai.u4aid24059@cb.students.amrita.edu}
\author[1]{\fnm{K} \sur{Nikhitha}}\email{cb.ai.u4aid24042@cb.students.amrita.edu}
\author[1]{\fnm{S Sai} \sur{Tejesh Reddy}}\email{cb.ai.u4aid24057@cb.students.amrita.edu}
\author*[1]{\fnm{Abhishek} \sur{S}}\email{s\_abhishek@cb.amrita.edu}

\affil*[1]{\orgdiv{Amrita School of Artificial Intelligence, Coimbatore}\\ \orgname{Amrita Vishwa Vidyapeetham}, \country{India}}

\abstract{This work presents an EEG-based classification framework for distinguishing motor actions and motor imagery using a consumer-grade Emotiv EPOC X headset. EEG recordings were collected from 47 participants using 14 channels sampled at 128~Hz while participants performed or imagined six limb movements involving the right and left arms, legs, and wrists. The resulting protocol formed twelve task classes by combining the body part with the action or imagery condition. The raw EDF recordings were processed using MNE-Python, including channel selection, 4--45~Hz band-pass filtering, signal-quality based trial rejection, per-recording normalisation and sliding-window segmentation. Each 10-second task segment was divided into 2-second windows with a 0.5-second stride, producing 8,092 windows from 476 trials. A heterogeneous representation of 1,771 features per window was then constructed using spectral, statistical, Hjorth, fractal and wavelet descriptors, together with spectrogram, scalogram, inter-hemispheric connectivity and Riemannian tangent-space features. An ensemble of Logistic Regression, Linear Discriminant Analysis, Extra Trees and Histogram Gradient Boosting was used as the main classifier, and Support Vector Machine, Random Forest, Naive Bayes, AdaBoost, EEGNet and CNN--LSTM models were evaluated for comparison. Subject-level GroupKFold validation was used to prevent windows from the same participant appearing in both training and evaluation folds. The ensemble achieved 80.25\% for binary action-versus-imagery classification, 26.68\% for six-class body-part classification and 21.43\% for the full twelve-class task, while left-versus-right imagery classification reached 52.28\% under nested cross-validation. The results show that combining complementary time--frequency, covariance-based and signal-derived features provides a strong baseline for subject-independent decoding of real versus imagined movement from low-cost EEG, while limb-level and side-level decoding remain challenging.}

\keywords{Electroencephalography (EEG), Brain--Computer Interface, Motor Imagery, Motor Action, Spectrogram, Scalogram, Riemannian Geometry, Machine Learning, Emotiv EPOC X}

\maketitle

%% ===================================================================
\section{Introduction}\label{sec1}

Brain--computer interfaces (BCIs) give the brain a channel to a computer that does not pass through muscles. They record brain activity, look for patterns that match an intention, and turn those patterns into a command. Motor applications are where this matters most. A person who has lost control of a limb after a stroke or a spinal injury can often still produce the cortical activity that would normally drive the movement, and a BCI that reads this activity could steer a wheelchair, open a prosthetic hand or guide a rehabilitation robot. Most such systems use electroencephalography (EEG). It is non-invasive and cheap compared with other recording methods, it follows brain activity on the scale of milliseconds, and it can be recorded with equipment small enough to carry \citep{Wolpaw2002,NicolasAlonso2012}.

Two related mental states are central to this work. In motor execution, or \emph{action}, the person really makes the movement. In motor imagery, the same movement is rehearsed in the mind while the body stays still. The two are not the same, but they share much of their neural machinery. Both reduce the power of the sensorimotor rhythms over the motor cortex, mostly in the mu band (8--12~Hz) and the beta band (13--30~Hz), and this drop, known as event-related desynchronisation, is the signal that most motor-imagery BCIs rely on \citep{Pfurtscheller2001,PfurtschellerAranibar1977}.

Much of the early work kept the problem small. The classic task is to tell imagined left-hand movement from imagined right-hand movement, and Common Spatial Pattern (CSP) filters followed by linear discriminant analysis or a support vector machine became the standard way to solve it \citep{Ramoser2000,Blankertz2008}. Adding more classes makes things far harder. The cortical areas for the hand, the wrist and the leg lie close together, the skull spreads their activity across the scalp, and the EEG of one person can look quite unlike that of the next \citep{Naeem2006,Lotte2018}. A model that works well for the people it was trained on can therefore fail on someone new.

The recording hardware adds its own limits. Clinical EEG amplifiers give clean signals, but they are expensive, slow to set up and rarely leave the laboratory. Consumer headsets such as the Emotiv EPOC are cheap, quick to put on and comfortable enough for everyday use, and studies have shown that they capture real neural activity \citep{Badcock2013,Duvinage2013}. They are also noisier, more sensitive to movement and limited in where they place their electrodes \citep{Maskeliunas2016}. The EPOC X used here has fourteen saline electrodes, none of which sits over the central motor strip (C3, Cz and C4), which is where motor activity is usually strongest.

Researchers have tried many ways of turning short EEG segments into something a classifier can use. Band-power and statistical features summarise each channel. Spectrograms and scalograms keep track of how the rhythms change over time. Riemannian methods describe a trial by the covariance between its channels and work directly with the geometry of these matrices \citep{Barachant2012,Congedo2017}. Deep networks such as EEGNet learn their own features from the raw signal \citep{Schirrmeister2017,Lawhern2018}. Each of these views captures something the others miss, so combining them is a natural step, particularly when the dataset is too small for a deep network to learn everything on its own.

How a model is evaluated matters as much as the model itself. EEG windows cut from the same recording are strongly correlated, so if windows of one person appear in both the training and the test set, the classifier can learn to recognise the person instead of the task. An early version of our own pipeline fell into exactly this trap and reported accuracies above 90\%, which dropped sharply once the evaluation was corrected. For a BCI that has to work for new users, the only fair test is one in which the people being tested were never seen during training.

In this study we built a complete classification pipeline around the Emotiv EPOC X and a twelve-task protocol. Participants performed and then imagined six movements: right-arm and left-arm flexion, right-leg and left-leg lifting, and right-wrist and left-wrist rotation. We analysed the same recordings at three levels of detail: action versus imagery (two classes), the body part involved (six classes), and the exact task (twelve classes). At our guide's suggestion we also asked whether the side of an imagined movement, left or right, can be decoded. Along the way we found that the task schedule did not start at a fixed time in every recording, and we recovered the true start from the headset's own motion sensor before cutting any trials. Every result we report comes from people the model had not seen during training.

The main contributions of this work are as follows:

\begin{enumerate}
\item \textbf{A new twelve-task EEG dataset of real and imagined limb movements.} We recorded 47 participants with a consumer-grade 14-channel Emotiv EPOC X while each of six movements (right and left arm, leg and wrist) was first performed and then imagined.

\item \textbf{Subject-independent decoding of executed versus imagined movement from consumer-grade EEG.} We show that real and imagined movement can be told apart in people the model has never seen, at 80.25\% with a 14-channel consumer headset. The result holds under person-wise cross-validation, leave-one-subject-out validation (80.25\%) and nested cross-validation (79.41\%).

\item \textbf{Classification of the same recordings at several levels of detail.} Every trial is classified in three ways: action or imagery (2 classes), the body part that was moved or imagined (6 classes), and the exact task, meaning body part and condition together (12 classes). In addition, an imagery-versus-imagery binary task uses only the imagery trials and asks whether the imagined movement was on the left or the right side of the body. Because all of these use the same trials, they show directly how accuracy falls as the question becomes more specific: 80.25\% for action versus imagery, 26.68\% for body part (chance 16.67\%), 21.43\% for the twelve tasks (chance 8.33\%), and 52.28\% for the imagery-versus-imagery (left versus right) task.

\item \textbf{A controlled comparison of feature representations and classifiers under one subject-independent protocol.} On identical folds we compare handcrafted, spectrogram, scalogram, connectivity, CSP and Riemannian features, and classical, ensemble and deep-learning models. The Morlet scalogram is consistently more informative than the STFT spectrogram, adding CSP features to each classifier changes accuracy by less than two points, and deep networks (EEGNet, CNN--LSTM) do not outperform feature-based models on a dataset of this size.
\end{enumerate}

The rest of the paper is organised as follows. Section~\ref{sec:related} reviews related work. Section~\ref{sec:methods} describes the protocol, data acquisition, preprocessing, feature extraction, classifiers and evaluation. Section~\ref{sec:results} presents and discusses the results, and Section~\ref{sec:conclusion} concludes the paper.

%% SECTIONS-BELOW
\section{Literature Review}\label{sec:related}

Decoding movement from scalp EEG has been studied for several decades, and the main obstacle has always been the same: the signal is weak, noisy and different from one person to the next. The physiological basis of almost all motor decoding is a change in the sensorimotor rhythms. When a person moves, or prepares to move, the power of the mu rhythm (8--12~Hz) and the beta rhythm (13--30~Hz) over the motor cortex falls, and it rises again, often above its resting level, once the movement ends. These two effects are called event-related desynchronisation and synchronisation (ERD/ERS) \citep{PfurtschellerAranibar1977,PfurtschellerLopes1999}. Neuper and Pfurtscheller showed that their frequency, timing and location depend on the body part involved, which is what makes it possible, in principle, to tell one movement from another \citep{Neuper2001}.

Motor imagery produces a similar pattern without any overt movement. Studies of the neural basis of imagery have found that imagining an action engages much of the same motor network as performing it, although usually more weakly \citep{Munzert2009}. This overlap is why imagery can drive a BCI, but it also means that action and imagery are not trivially separable. Pfurtscheller et al.\ found that imagined hand, foot and tongue movements each change the mu rhythm in their own characteristic way, and that these differences can be used to classify single trials, although not equally well for every person \citep{Pfurtscheller2006}. Most BCI work has concentrated on imagery alone. Distinguishing real from imagined movement, which matters for rehabilitation systems that must know whether a patient is trying to move or only thinking about it, has received far less attention.

Common Spatial Pattern (CSP) filtering is probably the most influential single method in the field. Ramoser et al.\ showed that spatial filters which maximise the difference in variance between two classes give strong results for imagined left- and right-hand movement \citep{Ramoser2000}. Blankertz et al.\ later explained why plain CSP is fragile when the data are noisy or the training set is small, and how regularised variants make it more stable \citep{Blankertz2008}. Filter-bank CSP (FBCSP) added the observation that the most useful frequency band differs from person to person, so it runs CSP in several sub-bands and keeps the most informative features \citep{Ang2012}. FBCSP won the four-class and two-class motor imagery tasks of BCI Competition IV and is still one of the standard baselines \citep{Tangermann2012}.

Covariance-based methods grew out of the limits of spatial filtering. Barachant et al.\ treated the covariance matrix of each trial as a point on the manifold of symmetric positive-definite matrices and classified trials by their Riemannian distance to the mean of each class \citep{Barachant2012}. Projecting the matrices onto a tangent space turns them into ordinary vectors that standard classifiers can use, while keeping the relationships between channels. Congedo et al.\ reviewed this family of methods and argued that they are both simpler and more robust than many earlier pipelines, especially for multi-class problems \citep{Congedo2017}.

Time--frequency analysis looks at the same signals from another angle. ERD/ERS unfolds over time, so averaging power over a whole window can hide part of it. The short-time Fourier transform (STFT) produces a spectrogram that shows how band power changes during a trial, and Tabar and Halici used STFT images of motor-imagery EEG as the input to a convolutional network \citep{Tabar2017}. The continuous wavelet transform (CWT) produces a scalogram instead. It analyses high frequencies with short windows and low frequencies with long ones, which suits signals whose rhythms come and go \citep{Torrence1998}. Simpler descriptors remain widely used as well: band power in fixed frequency bands \citep{Lotte2014}, the Hjorth parameters, which summarise the activity, mobility and complexity of a signal \citep{Hjorth1970}, and fractal measures of how irregular the signal is \citep{Higuchi1988}.

Deep learning is a further major direction. Schirrmeister et al.\ showed that convolutional networks trained on raw EEG can match FBCSP on motor imagery, and visualised what the networks had learned \citep{Schirrmeister2017}. Sakhavi et al.\ fed a temporal representation of filter-bank envelopes into a CNN to capture how the signal evolves within a trial \citep{Sakhavi2018}. EEGNet showed that a compact network built from depthwise and separable convolutions can work across several BCI paradigms with very few parameters \citep{Lawhern2018}. Reviews of deep learning for EEG, however, note that its advantage over well-designed classical pipelines is often small and uneven, particularly when the number of participants is limited \citep{Roy2019,Craik2019,Altaheri2023}.

Generalising to people who were not in the training set is still one of the hardest problems. The EEG of different people differs in amplitude, in the exact position of the electrodes and in how each person carries out the task, so a classifier trained on one group often performs poorly on another \citep{Lotte2018}. Alignment methods reduce this shift by re-centring each person's data before classification, either in Euclidean space \citep{He2020} or on the Riemannian manifold \citep{Zanini2018,Zhang2020}. The MOABB benchmark compared many algorithms on many public datasets and showed how strongly the ranking of methods can change from one dataset to the next, which argues for testing under realistic, subject-independent conditions rather than reporting a single favourable split \citep{Jayaram2018}.

Consumer-grade EEG lowers the cost of all this but adds problems of its own. Badcock et al.\ compared the Emotiv EPOC with a research system and found that it captured auditory event-related potentials of usable quality \citep{Badcock2013}. Duvinage et al.\ found it workable for a P300 speller, though clearly less accurate than a medical-grade amplifier \citep{Duvinage2013}, and Maskeli\={u}nas et al.\ reported high variability between trials when consumer headsets were used for control tasks \citep{Maskeliunas2016}. Motor-imagery BCIs have nevertheless reached clinical use; Ang et al., for example, combined CSP-based imagery detection with robotic feedback in a randomised trial of stroke rehabilitation \citep{Ang2015}.

Taken together, these studies leave a clear gap. Most motor-imagery work uses research-grade amplifiers, public benchmark datasets and two or four imagery classes. Far fewer studies separate real from imagined movement, compare several levels of detail on the same recordings, or test consumer hardware under strictly subject-independent evaluation. Our study addresses this gap. Table~\ref{tab:litreview} summarises the representative studies discussed above.

\begin{table*}[t]
\centering
\caption{Comparison of representative studies on EEG-based motor imagery and motor-related classification. N/R: not reported in a directly comparable form.}
\label{tab:litreview}
\footnotesize
\setlength{\tabcolsep}{2.5pt}
\begin{tabular}{@{}p{0.06\linewidth}p{0.15\linewidth}p{0.13\linewidth}p{0.17\linewidth}p{0.12\linewidth}p{0.12\linewidth}p{0.16\linewidth}@{}}
\toprule
Ref. & Task / Classes & Dataset / Hardware & Features / Representation & Model & Accuracy & Main Limitation \\
\midrule
\cite{Ramoser2000} & Left vs.\ right hand MI (2-class) & Research EEG, 3 subjects & CSP + log-variance & LDA & $\approx$80--90\% & Restricted to binary MI \\ \addlinespace
\cite{Pfurtscheller2006} & Hand, foot, tongue MI & Research EEG & Mu-band ERD/ERS patterns & Linear classifier & N/R & Imagery only; varies by person \\ \addlinespace
\cite{Naeem2006} & Hand, foot, tongue MI (4-class) & Research EEG & ICA-based representation & LDA & N/R & Limited spatial separability \\ \addlinespace
\cite{Ang2012} & Motor imagery (4-class) & BCI Competition IV-2a & Filter-bank CSP + feature selection & Naive Bayes & $\approx$68\% & Many frequency/filter parameters \\ \addlinespace
\cite{Barachant2012} & Multi-class MI & Research EEG & Covariance matrices + Riemannian geometry & MDM / tangent-space classifier & N/R & Costly covariance operations \\ \addlinespace
\cite{Tabar2017} & Motor imagery (2-class) & BCI Competition IV-2b & STFT images & CNN + stacked autoencoder & $\approx$77\% & Complex, data-sensitive training \\ \addlinespace
\cite{Schirrmeister2017} & Motor imagery (4-class) & BCI Competition IV-2a & Learned spatio-temporal features & Deep / Shallow ConvNet & $\approx$71--74\% & Data-hungry, less interpretable \\ \addlinespace
\cite{Sakhavi2018} & Motor imagery (4-class) & BCI Competition IV-2a & Temporal filter-bank representation & CNN & $\approx$74.5\% & More complex architecture \\ \addlinespace
\cite{Lawhern2018} & Several BCI paradigms & Public EEG datasets & Depthwise and separable convolutions & EEGNet & N/R & Needs enough training data \\ \addlinespace
\cite{He2020} & Cross-subject MI/ERP & Public EEG datasets & Euclidean alignment of covariance & Various classifiers & N/R & Assumes a moderate covariance shift \\ \addlinespace
\cite{Zanini2018} & Cross-subject MI & Public EEG datasets & Riemannian re-centring of covariances & MDM & N/R & Needs unlabelled data of the new subject \\ \addlinespace
\cite{Zhang2020} & Cross-subject EEG & Four EEG datasets & Riemannian alignment + tangent space & MEKT & N/R & Needs source/target adaptation \\ \addlinespace
\cite{Badcock2013} & EEG signal validation & Emotiv EPOC (14-ch) & Auditory event-related potentials & ERP analysis & N/R & Lower spatial resolution than clinical EEG \\ \addlinespace
\cite{Maskeliunas2016} & Concentration / blink control & Emotiv EPOC, NeuroSky & EEG control features & Conventional classifiers & N/R & High inter-trial variability \\ \addlinespace
\cite{Ang2015} & MI-based stroke rehabilitation & Stroke patients & CSP-based MI features & BCI + robotic feedback & N/R & Small clinical cohort \\ \addlinespace
\textbf{Our work} & Action vs.\ imagery (2), body part (6), twelve tasks (12), imagery vs.\ imagery left/right (2) & Emotiv EPOC X (14-ch, 128 Hz), 47 participants & 1,771 features: handcrafted + STFT + CWT + connectivity + Riemannian & Ensemble (LR + LDA + Extra Trees + HistGB), mirror augmentation, trial-level voting & 80.25\% (2-class); 26.68\% (6-class); 21.43\% (12-class); 52.28\% (imagery L/R, nested CV) & Consumer EEG with no central electrodes; limb and side remain hard to decode \\
\botrule
\end{tabular}
\end{table*}
\FloatBarrier

\section{Proposed Methodology}\label{sec:methods}
\emph{(to be added)}
\section{Results and Discussion}\label{sec:results}
\emph{(to be added)}
\section{Conclusion}\label{sec:conclusion}
\emph{(to be added)}

\bibliography{bibliography}

\end{document}
